\chapter{Fusion, вывод и консолидация}
\label{ch:fusion}

Обход достаёт то, что уже есть в графе. Но память живого существа ещё и
порождает новое: связывает далёкое, делает выводы, обобщает. Этот класс
операций мы называем fusion. Здесь разбирается, как именно из активированных
узлов рождается новый узел, чем вывод отличается от ассоциации и обобщения и
как всё это остаётся под эпистемическим контролем.

\section{Что такое fusion, механически}

\begin{definition}[Fusion, подробно]
Fusion --- оператор, который по совместно активированной конфигурации узлов
$\{A_1, \ldots, A_n\}$ создаёт новый узел $X$ и связывает его с посылками
рёбрами провенанса (обоснования). Новый узел рождается провизорно: с низкой
уверенностью и, при необходимости, с флагом \texttt{need\_to\_discover} на том,
что его различает.
\end{definition}

Механика в четырёх шагах:
\begin{enumerate}
\item Со-активация. Обход (retrieval, любопытство или сон) подсвечивает
несколько узлов одновременно.
\item Распознавание конфигурации. Оператор замечает, что активная конфигурация
согласована --- образует осмысленный паттерн (причинную цепочку, аналогию,
повтор).
\item Рождение узла. Создаётся $X$ с содержанием (сформулированным языковой
моделью) и провенансом $\{A_1, \ldots, A_n\}$.
\item Провизорная запись и пропагация. $X$ пишется с малым \texttt{conf},
затем пропагируется со знаком, карантином и дисконтом независимости
(глава~\ref{ch:propagation}).
\end{enumerate}

\begin{principle}[Fusion рождает гипотезу, а не факт]
Результат fusion --- всегда кандидат, а не запись. Это защищает от превращения
обаятельной ассоциативности в фабрику ложных убеждений
(глава~\ref{ch:epistemics}). Повышение до «факта» --- только через накопление
подтверждений или выживание в ночном скрабе.
\end{principle}

\section{Три вида fusion}

\smallskip\noindent\textbf{Ассоциация.}
Два далёких узла со-активируются, и рождается связь (или узел-мостик). Это
источник «неожиданных, но осмысленных» ходов мысли.

\begin{workedbox}[Ассоциация: Платон-боец]
Ночной обходчик случайно со-активирует два узла: «Платон был олимпийским
чемпионом по панкратиону (боец/борец --- по свидетельству Олимпиодора)» и
«\textit{platon} по-гречески --- „широкий, широкоплечий“». Fusion замечает
согласование «мощное телосложение бойца» + «имя означает „широкий“» и
рождает гипотезу: «„Платон“ --- прозвище по телосложению, а не данное имя»,
\texttt{conf = 0,2}, провенанс из двух узлов, флаг \texttt{need\_to\_discover}
на «проверить: прозвище ли это и каково настоящее имя». Вывод правдоподобен
(версий, к слову, несколько, но все сводятся к трактовке слова «широкий», а
настоящее имя по преданию известно) --- но именно гипотеза: если не пометить
низкой уверенностью, это тот самый преждевременный коллапс
(глава~\ref{ch:epistemics}). Существенно, что связывает не «Платон-философ»
сам по себе, а именно пара «боец» + «перевод».
\end{workedbox}

\smallskip\noindent\textbf{Вывод (инференция).}
Из посылок рождается заключение. Здесь принципиально важен случай $n$-арного
вывода, где ни одна часть посылок не достаточна.

\begin{workedbox}[$N$-арный вывод: корм котам]
Пусть в памяти Рин про пользователя есть три узла, и каждый подтягивает свой
контекст из эмбрионального графа (см. раздел ниже):
\begin{itemize}
\item $A$ --- «у него есть коты» $\Rightarrow$ контекст: котов надо кормить,
нужен корм;
\item $B$ --- «живёт в Португалии» $\Rightarrow$ контекст: магазины закрываются
рано;
\item $C$ --- «иногда ездит на пары» $\Rightarrow$ контекст: с пар возвращается
поздно.
\end{itemize}

Только совмещение всех трёх с их контекстами даёт вывод:
\[
\begin{aligned}
X =\ &\text{«купи котам корм на выходных --- в будни по пути с пар} \\
     &\text{можешь не успеть до закрытия»}.
\end{aligned}
\]

Существенно, что ни одна пара не достаточна: без «магазины рано» нет спешки,
без «поздно с пар» нет проблемы, без «нужен корм» нет повода. Вывод рождается
лишь при одновременном присутствии трёх посылок и их контекстов.
\end{workedbox}

Как мы это представляем сейчас, не меняя структуру: fusion создаёт узел $X$ с
провенансом $\{A,B,C\}$. Тем самым тройка-посылка кодируется как
узел-фактор, ссылающийся на три посылки, --- это factor-graph-кодировка
гиперребра $\{A,B,C\} \to X$. Иначе говоря, наша память уже наполовину
гиперграф: попарные рёбра $A$--$B$, $B$--$C$ могут отсутствовать (Португалия и
коты не связаны напрямую), а $n$-арный вывод всё равно рождается --- через
общий узел-обоснование. Стоит ли делать гиперрёбра первоклассными объектами
(со своим обходом и лапласианом) --- тема приложения~\ref{ch:hypergraphs}.

\begin{figure}[h]
\centering
\begin{tikzpicture}[
    node distance=1.3cm and 2.0cm,
    box/.style={draw, rounded corners, align=center, inner sep=6pt},
    prov/.style={draw=none, align=center}
]
\node[box] (a) at (0,1.2) {$A$: Португалия};
\node[box] (b) at (0,0) {$B$: коты};
\node[box] (c) at (0,-1.2) {$C$: универ};
\node[box] (x) at (5,0) {$X$: купи корм\\на выходных};
\node[prov] (p) at (2.5,1.2) {provenance};
\draw[-{Latex[length=2mm]}] (a) -- (x);
\draw[-{Latex[length=2mm]}] (b) -- (x);
\draw[-{Latex[length=2mm]}] (c) -- (x);
\end{tikzpicture}
\caption{$N$-арный вывод как узел-фактор $X$ с провенансом $\{A,B,C\}$.
Попарных рёбер между $A$, $B$, $C$ может не быть.}
\end{figure}

\smallskip\noindent\textbf{Как посылка «подтягивает свой контекст».}
Важная деталь, без которой $n$-арный вывод не работает: посылки участвуют не
«голыми», а вместе со своим контекстным ореолом. Механически это ограниченная
(1--2 хопа), смещённая целью/настроением spreading activation от каждой посылки
в эмбриональный граф:
\[
A = \text{коты} \Rightarrow \{\text{нужен корм}\},
\]
\[
B = \text{Португалия} \Rightarrow \{\text{магазины рано закрываются}\},
\]
\[
C = \text{универ} \Rightarrow \{\text{поздний возврат}\}.
\]

Fusion работает над объединением этих ореолов, а не над тремя голыми узлами;
композитор (языковая модель) собирает из «нужен корм + магазины рано + поздний
возврат» связный actionable-вывод. Провенанс $X$ --- посылки и реально
использованные контекст-узлы. Ореол смещён текущей целью: иначе подтянулась бы
вся Вселенная соседей, поэтому активируется лишь релевантное намерению (тот же
mood/goal-bias, что в retrieval, глава~\ref{ch:walker}). Именно поэтому
эмбриональный граф (глава~\ref{ch:embryonic}) не декорация: без рёбер
ConceptNet/ATOMIC у посылок нет ореолов --- и вывод не рождается.

\begin{openq}[Трудная часть --- распознать связность («трудные 80\%»)]
Подтянуть ореолы дёшево; трудно распознать, что их объединение образует
связную цепочку рассуждения, а не случайную кучу со-активаций. Кандидаты
разобраны в главе~\ref{ch:timescales} (раздел~\ref{sec:whatfusiondoes}):
LLM-судья / графовые шаблоны-схемы / обучаемый скорер связности --- ни один
пока не признан удачным. На старте --- LLM-композитор предлагает, а
$\varphi$-скорер и обратная связь фильтруют
(главы~\ref{ch:personality}, \ref{ch:voice}). Это открытый инженерный вопрос.
\end{openq}

\smallskip\noindent\textbf{Обобщение (абстракция).}
Много однотипных эпизодов сливаются в один узел более высокого уровня. Это и
есть консолидация (следующий раздел).

\section{Fusion сцепляется: sub-fusion --- это тоже fusion}

Виды fusion --- не разные механизмы, а один оператор, порождающий узлы разного
рода. И fusion сцепляется: обобщение, построенное одним шагом («sub-fusion»),
становится посылкой для вывода на следующем. Абстракция поднимает уровень,
вывод делает конкретный скачок, опираясь на неё:
\[
\text{факты} \xrightarrow{\text{абстракция}} \text{обобщение}
\xrightarrow{\text{вывод}} \text{гипотеза}.
\]

Механика та же (со-активация $\to$ новый узел с провенансом); «sub» лишь
означает, что этот узел --- промежуточное обобщение, мост для последующего
шага. (Это, к слову, клеточная картинка из приложения~\ref{ch:cellcomplex}:
высшая клетка, приклеенная к низшим.)

\begin{workedbox}[Дневная находка и полная эпистемическая петля: House / Cast Away]
\textbf{Фон.} Фоновый разведчик (глава~\ref{sec:background}), даже на скудном
дневном бюджете, со-активирует два далёких узла: «в House M.D. есть Уилсон» и
«в Cast Away есть „Уилсон“ --- волейбольный мяч».

\textbf{Sub-fusion (абстракция).} Оба узла порождают обобщение:
«Уилсон --- приём экстернализации внутреннего монолога ГГ» (у Cast Away --- из-за
одиночества на острове; у House --- потому что ГГ умнее всех, детектив,
интроверт). Родился узел более высокого уровня --- общая сценарная функция.

\textbf{Fusion (скачок).} Эта абстракция + совпадение имени «Уилсон» рождают
гипотезу: «а вдруг имя Уилсон в House --- отсылка к Cast Away?!» Провизорно,
низкая уверенность.

\textbf{Проверка (петля эпистемики, глава~\ref{ch:epistemics}).}
\texttt{need\_to\_discover} на «годы выхода» $\to$ внешняя проверка: Cast Away
(2000) раньше пилота House (2004) --- совместимо, уверенность растёт. Затем
\texttt{need\_to\_discover} на «есть ли ещё свидетельства» $\to$ проверка
опровергает: официальная версия --- оммаж Шерлоку Холмсу (House $\approx$
Holmes, адрес 221B), а Джеймс Уилсон $\approx$ доктор Ватсон (одна буква,
созвучие). Гипотеза понижается, на её место встаёт лучше обоснованная.

\textbf{Мораль.} Красивый, согласующийся с частью данных вывод оказался
ложным --- это тот самый преждевременный коллапс, пойманный проверкой. И
заметьте: «правильный» ответ (Holmes $\to$ House) --- сам по себе fusion,
просто с лучшим провенансом. Как говорит ГГ, «все лгут» --- поэтому
уверенность держим ниже единицы и идём проверять.
\end{workedbox}

\section{Консолидация как огрубление}

\begin{definition}[Консолидация, подробно]
Консолидация --- оператор, заменяющий множество похожих узлов/рёбер одним
обобщённым, сохраняя релевантные свойства и отбрасывая детали. Формально это
огрубление графа (coarse-to-local, глава~\ref{ch:timescales}).
\end{definition}

\begin{workedbox}[Консолидация: «Z --- частый гость»]
Десять эпизодов «Z был на стриме в понедельник», «...в среду», ... сворачиваются
в один узел «Z --- частый гость, любит доту». Выживают: частота, тёплая
валентность, тема (дота). Теряются: точные даты и реплики. Плохое огрубление
сохранило бы даты и потеряло валентность --- Рин помнила бы, когда заходил $Z$,
но забыла бы, как к нему относится. Значит критерий консолидации --- сохранять
аффективно- и реляционно-релевантное, а не всё подряд.
\end{workedbox}

Консолидация запускается двумя путями: во сне (фоновая переработка,
глава~\ref{ch:sleep}) и забыванием (когда эпизоды тускнеют, их дешевле слить в
обобщение, чем хранить порознь, глава~\ref{ch:timescales}).

\section{Кто пишет в узел и в каком виде}
\label{sec:whowrites}

Отдельно разберём вопрос, который легко замять: кто именно порождает содержание
нового узла --- LLM, автомат или Rinness --- и в каком виде оно ложится в
память. Полезно разбить на три подшага, потому что механизмы у них разные.

\smallskip\noindent\textbf{Решение (сливать ли).}
Для fusion это распознавание связной конфигурации (открытый «трудный 80\%»,
раздел~\ref{sec:whatfusiondoes}); для консолидации --- обнаружение избыточного
кластера. И тут важная поправка: группировать «однотипное» по строке нельзя.
Зритель мог «заходить», а мог «забегать» --- чистый автомат, ключующийся по
глаголу, тут же ломается, ведь мы не типизируем глаголы, которыми к Рин
обращаются. Синонимию ловит не автомат, а либо LLM при записи (нормализует
«забежал» в канонический тип связи \texttt{visited}), либо кластеризация по
близости эмбеддингов при консолидации. Поэтому честная картина: автомат делает
бухгалтерию (счётчики, хаб, перекидка рёбер, усреднение валентности),
эмбеддинги --- группировку «похоже по смыслу», а маленький LLM --- итоговый
ярлык. «Только автоматом» проходит лишь случай точных дублей.

\smallskip\noindent\textbf{Содержание узла.}
Новое осмысленное содержание (вывод, абстракция, ярлык обобщения) пишет
языковая модель --- шаблоном такое не собрать. Но та же самая модель, а не
какая-то отдельная: «нейтральный режим» и «Rinness» --- это два настроения
одной модели, а не два мозга. Для записи в память её держат в спокойном,
откалиброванном «девелоперском» режиме вектора настроения: рассудочно, минимум
стиля. Если голос Рин зашит в веса (LoRA), совсем нейтральной её не сделать ---
останется лёгкий налёт, и это не страшно, лишь бы не перегибать: в память важно
содержание, а тон всё равно родится заново на выходе. Автомат при этом ставит
все метаданные (провенанс, уверенность, \texttt{need\_to\_discover},
\texttt{importance}, домен) --- детерминированно.

\smallskip\noindent\textbf{Кто именно пишет --- и почему это не «Rinness».}
Пишет в граф та же Рин, но в спокойном режиме записи, а не в режиме голоса.
Здесь важно не путать «Рин» и «Rinness»: Rinness --- это узко выходная
стилизация (голос, гл.~\ref{ch:voice}), которая превращает извлечённое
содержание в реплику пользователю. Вот эта стилизация в граф не пишет. Одна
Рин, два режима: режим записи (спокойный) кладёт нейтральное содержание, режим
голоса (Rinness) звучит на выходе. Отсюда принцип: храни спокойно, стилизуй на
выходе. Причина простая: узел переиспользуется в разных будущих настроениях, и
если вморозить в него сегодняшний тон, он будет отдавать им и через месяц,
когда Рин к тому же зрителю потеплеет.

\smallskip\noindent\textbf{В каком виде хранить: содержание против стиля.}
Итак, в память кладём нейтральное содержание («зритель Z проиграл в доте») плюс
её отношение в метаданных --- валентность и, для значимых случаев, отдельный
узел-реакцию. Личность памяти живёт в отношении, а не в замороженных словах.
Класть готовую фразу «в голосе» («Z снова слил, ну куда ему») --- плохо: это
вмораживает транзиентный тон, память пухнет, а стилизовать на выходе уже
нечем.

Исключение --- вспышечная память. У людей рядовые события помнятся по сути, а
редкие пиково-эмоциональные --- вплоть до того, что было сказано и как
ощущалось. По аналогии: для редких моментов с высокой значимостью (большой
$|\Delta \mood|$) можно дополнительно сохранить дословный снимок --- реальную
реплику/тон того мгновения --- как отдельный связанный узел, а не как
каноническое содержание. Так по умолчанию память нейтральна и переиспользуема,
но яркие пики сохраняют свою фактуру.

\section{Выход fusion --- гипотеза, а не установленный факт}

Снимем терминологическую путаницу, потому что от неё зависит понимание всего.
Каждый узел --- это убеждение с уверенностью $\conf \in [-1,1]$. «Факт» в
широком смысле --- это просто любой узел, и да, fusion создаёт узел. Но узлы
различаются по уверенности: прямое наблюдение или измерение получает высокую
уверенность, а вывод (fusion) --- низкую. Именно это имелось в виду под
«fusion не коммитит факты»: коммит --- это запись как установленного,
высокоуверенного убеждения, которым дальше пользуются как истиной; fusion же
рождает узел-гипотезу. Повышение гипотезы до установленного статуса происходит
позже --- через корроборацию или проверку, а не в момент рождения.

Почему это критично: без такого разделения обаятельная ассоциативность
превратилась бы в фабрику ложных убеждений (пример House / Cast Away выше).
Поэтому fusion смел, но не безрассуден: гипотезы рождаются свободно, но влияние
гипотезы на соседей гейтируется её собственной (низкой) уверенностью и
надёжностью посылок --- карантин, глава~\ref{ch:propagation}; а различающие
пробелы, которых Рин не знает, помечаются \texttt{need\_to\_discover}, чтобы
любопытство или внешний вход их потом добыли (глава~\ref{ch:epistemics}).
Дальнейшую судьбу гипотезы --- проверку, повышение или понижение --- ведёт уже
эпистемический слой; fusion лишь честно рождает её с низкой уверенностью и
передаёт дальше.

\begin{principle}[На первых версиях: генератор кандидатов с аудитом]
Fusion --- главный источник и магии, и мусора. Пока распознавание связности
со-активации не решено (открытый вопрос), fusion на ранних версиях должен быть
не «автоматической когнитивной магией», а генератором кандидатов с аудитом.
Дефолтная позиция смещена в сторону осторожности: больше кандидатов в
\texttt{pending}, меньше \texttt{commit}. Принятые fusion потом ревизуются
отложенно (precision после проверки, глава~\ref{ch:experiment}), и только
выжившие повышаются. Лучше недодать смелых узлов, чем зарастить память красивым
мусором; смелость безопасно включим, когда научимся мерить связность.
\end{principle}

\begin{principle}[Разведи candidate\_text и canonical\_node]
Fusion всё ещё сильно опирается на LLM-композитор, и это риск: модель выдаёт
красивый текст, но плохо индексируемый и плохо ревизуемый узел. Разделяем два
формата. LLM формулирует \texttt{candidate\_text} --- гипотезу человеческим
языком. Но коммит идёт не текстом, а через schema-writer в \texttt{canonical\_node}
со структурой:
\texttt{\{node\_kind, claim\_type, entities, time\_scope, confidence, provenance, distinguishing\_unknowns\}}.
Иначе получишь хорошо написанные, но нетипизированные узлы, которые нельзя ни
точно найти (forensic-ретрив, гл.~\ref{ch:runtime}), ни аккуратно ревизовать.
Красота --- на входе, схема --- на коммите.
\end{principle}

\section{Когда со-активация --- повод для реплики, а не для узла-факта}
\label{sec:antifusion-joke}

Не всякая яркая со-активация должна оседать убеждением. Иногда две всплывшие
вещи --- это повод что-то сказать, а не повод что-то запомнить как факт о мире.
Это ключевое различение: полезное для реплики $\neq$ достойное записи фактом
(anti-fusion). Разберём на каноничном случае.

\begin{workedbox}[Чапаев и «нюанс»: шутку не пишем, реакцию --- пишем]
В памяти со-активировались «Чапаев» и слово «нюанс». Walker сводит их рядом, и
всплывает ассоциация: «тут просится анекдот». Что происходит по нашим
механизмам.

(1) \textbf{Откуда берётся анекдот.} Не из графа как хранимый текст. Граф даёт
только триггер --- «эта пара уместна для шутки». Сам текст анекдота порождает
базовая LLM (она знает его из предобучения) --- ровно как с обучением: граф
всплыл повод, модель поставляет содержание.

(2) \textbf{Что НЕ записывается.} Сам анекдот не коммитится узлом-фактом. Связь
«Чапаев relates-to нюанс через этот анекдот» --- это \texttt{candidate\_text}
для реплики, а не убеждение о мире; anti-fusion-гейт его как факт не
пропускает. Иначе память заросла бы «фактами», которые на деле каламбуры.

(3) \textbf{Что записывается.} Два узла другого рода. Во-первых, episode:
«рассказала анекдот про Чапаева-нюанс собеседнику $Z$» (со временем и
провенансом) --- это было, это факт о событии, не о мире. Во-вторых, реакция
$Z$ обновляет \texttt{reaction\_pattern} через Beta-постериор, условленный на
отношении (гл.~\ref{ch:personality}): зашло $\to \alpha{+}1$, поморщился
$\to \beta{+}1$. Salience-гейт решает судьбу: сдвинула реакция настроение
($|\Delta \mood|$) --- консолидируется; «ну ок» --- тускнеет.

(4) \textbf{Бонус --- крючок, а не шутка.} Если такой юмор у $Z$ стабильно
заходит, осесть может лёгкий \texttt{reaction\_pattern}: «пара
Чапаев+нюанс $\to$ этот анекдот заходит $Z$». Заметьте: выучивается модель того,
что кому заходит, а не сам анекдот хранится фактом. Комедийный тайминг ---
Beta-апостериор, не строка в базе.
\end{workedbox}

Мораль в одну строку: fusion-подобная со-активация может кончиться репликой, а
осесть в памяти должны событие речи и реакция на неё --- потому что важна не
шутка, а то, как она легла. Тот же принцип защищает память от превращения в
свалку острот, мемов и подходящих-к-случаю ассоциаций.

\begin{intuition}[Инженер, а не фантазёр]
Правильная Рин --- «инженер, которую делают бывшие физики-фундаменталисты»:
она охотно строит гипотезы, но по умолчанию ставит уверенность ниже единицы и
идёт проверять. Fusion даёт первое (смелость ассоциаций), эпистемическая
гигиена (глава~\ref{ch:epistemics}) --- второе (дисциплину). Одно без другого
--- либо сухой поисковик, либо обаятельный врун.
\end{intuition}
